% TeXlate 评测体系与指标演化全景报告
% 编译: latexmk -xelatex report.tex  (或 xelatex ×2)
\documentclass[11pt,a4paper]{ctexart}

\usepackage{geometry}
\geometry{margin=2.4cm,top=2.6cm,bottom=2.6cm}
\usepackage{pgfplots}
\pgfplotsset{compat=1.18}
\usepgfplotslibrary{dateplot,groupplots,statistics}
\usepackage{tikz}
\usetikzlibrary{positioning,arrows.meta,shapes.geometric,fit,backgrounds,calc}
\usepackage{booktabs}
\usepackage{multirow}
\usepackage{caption}
\captionsetup{font=small,labelfont=bf}
\usepackage{xcolor}
\usepackage{hyperref}
\hypersetup{colorlinks=true,linkcolor=blue!60!black,urlcolor=blue!60!black,citecolor=blue!60!black}
\usepackage{enumitem}
\setlist{nosep,leftmargin=1.6em}
\usepackage{fancyhdr}
\pagestyle{fancy}
\fancyhf{}
\fancyhead[L]{\small TeXlate 评测体系与指标演化}
\fancyhead[R]{\small 2026-09-19}
\fancyfoot[C]{\small\thepage}
\renewcommand{\headrulewidth}{0.4pt}

% 图配色
\definecolor{cA}{HTML}{1f6fb2}   % 主蓝
\definecolor{cB}{HTML}{c0392b}   % 红
\definecolor{cC}{HTML}{8a6d00}   % 棕
\definecolor{cD}{HTML}{2e8b57}   % 绿
\definecolor{cE}{HTML}{7d3c98}   % 紫
\definecolor{cF}{HTML}{16a085}   % 青

\title{\bfseries TeXlate 评测体系与指标演化全景报告\\[4pt]
  \large 七臂 benchmark · corpus\_v3 全量综合测试 · 2026-09-19}
\author{texlate bench 工作组}
\date{2026 年 9 月 19 日}

\begin{document}
\maketitle

\begin{abstract}
\noindent 本文是 TeXlate 评测体系的全景报告：上半部分为不熟悉项目的读者提供背景（管线架构、为什么 arXiv LaTeX 翻译难、benchmark 设计），下半部分是数据——全部历史 commit 考古出的指标演化时间线、2026-09-19 在 corpus\_v3 全量 13{,}266 篇上的综合测试画像、失败归因图谱与已知缺口。全部数字可追溯：每个指标注明口径与结果目录，数据与图源文件随文档入库。
\end{abstract}

\vspace{1em}
\tableofcontents
\newpage

\section{背景知识}

\subsection{这个项目做什么}

TeXlate 是「幻觉翻译」(hjfy.top) 的开源复刻：输入一篇 arXiv 论文的 LaTeX 源码（e-print），输出一份重编译的中文 PDF，正文逐段替换为中文译文、数学公式/图表/参考文献结构原样保留，并提供双语对照阅读。链路上的每一段都由自研组件构成：

\begin{itemize}
  \item \textbf{获取层}：arXiv e-print 批量下载（IA bulk / export / OAI-PMH 三通道），钉版入库；
  \item \textbf{解析层}：自研半解析器——不是完整 TeX 引擎，而是「逐字符扫描 $\to$ pieces $\to$ 占位符保护 $\to$ 区间 splice」的保真切分器，外加宏展开机（Gullet）把用户自定义宏就地展开；
  \item \textbf{翻译层}：段落级 chunk 经 LLM 网关批量翻译（swe-2 系列模型），占位符协议保护公式与命令不被译文污染；
  \item \textbf{校验层}：L0/L1/L2 三层校验（chunk 不变量、tree-sitter 结构、LLM judge）；
  \item \textbf{编译层}：ctex 注入 + normalize + xelatex/tectonic 双引擎编译；
  \item \textbf{修复层}：fixloop——131 条 YAML 规则驱动的自动修复循环（装包、改源、路由、polyfill），把编译失败的论文迭代救回；
  \item \textbf{对齐层}：named-destination 锚点同步，滚动阅读时中英逐段对齐。
\end{itemize}

\subsection{为什么这件事难}

\textbf{arXiv 源码是脏的真实世界数据}。语料横跨 LaTeX2.09（1990s）到 2026 年：混用 plain \TeX{} 原语与 LaTeX2.09 旧式、非 UTF-8 编码（Big5/GBK/Latin-1）、缺文件的残缺包、tar/gz 双格式、古早宏包（revtex4、aastex、psfig）。基线实测：裸编译（无修复）union pdf 仅 65--72\%，clean 22--40\%——三分之二的论文源码直接编译不出 PDF，最大失败类是 missing\_file（占裸失败 92.1\%）。

\textbf{宏展开是现成解析器的共同盲区}。95\% 的论文含自定义宏（中位 47 个/篇，最多 446），其中 12/39 篇手挑语料存在「宏体里藏 \verb|\begin/\end|」的结构性陷阱，49\% 的宏体藏 $>20$ 字母的可译文本。对 8 个第三方解析库（latex-utensils、unified-latex、tree-sitter-latex、TexSoup、plasTeX、pylatexenc、LaTeX.js、ieeA）的横评（\texttt{bench/results/00-grand-comparison.md}）表明：陷阱断言全军覆没——最好成绩 24/26 vs 自研半解析器 32/32；「能 parse」不等于「parse 对」，pylatexenc 报 90/90 成功却是静默截断零报错。\textbf{无错误信号比崩溃更危险}，这决定了主解析器必须自研、校验器必须独立。

\textbf{翻译约束是结构性的}。chunk 边界纪律要求公式、命令、环境标签逐字节保留，译文只能在散文区间替换；术语表一致性、学术语域（register）、不可译段（参考文献/版权行）都要显式协议。编译侧中文需要 ctex 注入与 CJK 字体路由，修复侧要处理 EPS 图源、缺 sty/cls、bib 链路断点等 131 类已建档机制。

\subsection{管线架构}

图~\ref{fig:pipeline} 给出全链架构与各段的 benchmark 对应关系。

\begin{figure}[htbp]
\centering
\resizebox{\textwidth}{!}{%
\begin{tikzpicture}[
  font=\small,
  stage/.style={draw,rounded corners=2pt,minimum height=8mm,inner xsep=4pt,align=center,fill=cA!8},
  impl/.style={font=\scriptsize\color{black!60},align=center},
  bench/.style={draw=cC,dashed,rounded corners=1pt,font=\scriptsize\color{cC},inner sep=2pt,align=center,fill=cC!6},
  arr/.style={-{Stealth[length=2mm]},thick,black!70},
  node distance=5mm and 4mm
]
% 主流水线
\node[stage] (fetch) {arXiv\\e-print};
\node[stage,right=of fetch] (parse) {flatten + 半解析\\mouth/gullet/segmenter};
\node[stage,right=of parse] (xlat) {段落级翻译\\xlat 编排 + 网关};
\node[stage,right=of xlat] (valid) {校验\\L0/L1/L2};
\node[stage,right=of valid] (splice) {splice 重建\\+ ctex 注入};
\node[stage,right=of splice] (comp) {编译\\xelatex / tectonic};
\node[stage,right=of comp] (align) {锚点对齐\\+ 双语 PDF};
\draw[arr] (fetch) -- (parse);
\draw[arr] (parse) -- (xlat);
\draw[arr] (xlat) -- (valid);
\draw[arr] (valid) -- (splice);
\draw[arr] (splice) -- (comp);
\draw[arr] (comp) -- (align);
% fixloop 回环
\node[stage,below=6mm of comp,fill=cB!10,draw=cB] (fix) {fixloop 修复循环\\131 条 YAML 规则};
\draw[arr,cB] (comp.south) -- (fix.north);
\draw[arr,cB] (fix.west) -| (splice.south);
% bench 标注
\node[bench,below=4mm of parse] (b1) {B1 parsebench\\identity/leak/flatten};
\node[bench,below=4mm of xlat] (b4) {B4 xlatbench + qualbench\\硬契约/延迟/ESA 评分};
\node[bench,below=4mm of valid] (b6) {B6 validbench\\检出率/FP};
\node[bench,below=21mm of splice] (b3) {B3 compilebench\\clean/pdf 率};
\node[bench,below=4mm of align] (b7) {B7 alignbench\\named-dest 保留};
\node[bench,above=4mm of xlat] (b5) {B5 e2ebench 全链漏斗};
\draw[cC,dashed,-{Stealth[length=1.5mm]}] (b5.south) -- (xlat.north);
\node[bench,above=4mm of parse] (b2) {B2 fixtures T01--T98};
\end{tikzpicture}}%
\caption{TeXlate 管线架构与 benchmark 臂对应关系。主流水线自左向右；fixloop 是编译段的自动修复回环；B1--B7 为七个评测臂（虚线框），另有 fixtures 陷阱断言集做解析段单元级测试。}
\label{fig:pipeline}
\end{figure}

\subsection{术语与口径}

\begin{itemize}
  \item \textbf{cell / 格}：语料库中一篇论文的存储单元（\texttt{\{id\}/\{meta.json, raw.*, extracted/\}}）；benchmark 语境下也指一个被评对象（paper $\times$ engine 组合）。
  \item \textbf{层 / layer}：corpus\_v3 的抽样分层（见 \S\ref{sec:corpus}），每层独立 manifest、独立抽样框。
  \item \textbf{臂 / arm}：同一被测对象的不同处理条件——compile 的 baseline（裸编译）vs zh（产品链 normalize 后）；e2e 的 mock（假翻译）vs real（真 LLM）；stagerun 的 base vs zh 条件。
  \item \textbf{clean / pdf\~ / FAIL}：编译终态分级——clean 无 warning 出 PDF；pdf\~ 带可接受瑕疵出 PDF；FAIL 无 PDF。union pdf = clean + pdf\~。
  \item \textbf{identity / leak}：解析质量两轴——identity 为 reconstruct 逐字节一致率；leak 为受保护内容（公式/命令）漏进可译 chunk 的占比。
  \item \textbf{scorecard}：gate\_scorecard.py 对 stagerun records 的终态统计，M2 出口门为 union pdf $\geq 90\%$、clean $\geq 90\%$。
  \item \textbf{机制 / mechanism}：mechanisms.jsonl 台账登记的一类已知缺陷模式（如 dollar 族泄漏、eps\_route），benchmark 发现 $\to$ 归因 $\to$ 规则沉淀的反馈环。
\end{itemize}
   % 背景知识 + 管线架构
\section{评测体系设计}

\subsection{七臂总览}

benchmark 套件按「每臂对准管线一段 + 一个组合层」设计，评测器规格冻结件随档于 \texttt{refs/}（\texttt{10-benchmark-suite.md}），共同底材为 corpus\_v3。表~\ref{tab:arms} 为总览。

\begin{table}[htbp]
\centering\small
\caption{七臂 benchmark 总览（B1--B7 + 辅助件）}
\label{tab:arms}
\begin{tabular}{@{}lllll@{}}
\toprule
\# & benchmark & 测哪段 & 底材 & 核心指标 \\
\midrule
B1 & parsebench & 解析段 & corpus\_v3 + corpus39 & ok / identity / leak / dead / flatten \\
B2 & fixtures & 解析段（单元级） & 手构陷阱 \texttt{.tex} 集（10 件） & 陷阱断言通过率（T01--T98） \\
B3 & compilebench & 编译段 + fixloop & corpus\_v3 raw blob & clean/pdf\~{}/FAIL、救回率、规则谱 \\
B4 & xlatbench & 翻译段 & corpus\_v3 chunk 抽样 & 硬契约率 / 延迟 / token 经济 \\
B4b & qualbench & 翻译质量 & e2e state 配对 & LLM-judge ESA 评分 / flag 分类 \\
B5 & e2ebench & 全链（组合） & corpus\_v3 子集 & 逐段漏斗 + 终态分布 + 引入回归 \\
B6 & validbench & 校验段 & 合成破坏案例 & 检出率 / FP / 延迟 \\
B7 & alignbench & 阅读体验 & en/zh 编译产物对 & named-dest 保留率 / 链权 \\
\midrule
\multicolumn{5}{@{}l}{辅助件：stagerun（五阶段批量 DAG）、gullet\_bench（展开机实测）、}\\
\multicolumn{5}{@{}l}{gate\_scorecard（M2 出口门）、translators\_bench（mock 臂工厂）、nightwatch}\\
\bottomrule
\end{tabular}
\end{table}

依赖序：corpus\_v3 $\to$ B1 $\to$ B4 $\to$ \{B3, B5\} $\to$ B7；B2/B6 以合成输入独立。

\subsection{评测现场：\texttt{bench/} 的构成}

所有评测资产集中在仓库 \texttt{bench/} 树，入口文档为 \texttt{PROTOCOL.md}（逐库横评协议）与 \texttt{TIERS.md}（验证分层契约，均随档 \texttt{refs/}）。后者定义「一个问题该在哪一级被抓住」：L0 单测/fixture 断言（秒级，CI 硬门）$\to$ L1 corpus 机制覆盖（分钟级）$\to$ L2 stagerun 子集回归（分钟–小时，机制台账定向重放）$\to$ L3 全层全臂过夜集成（里程碑门）；规则是「能低不高」——低层能钉的不许只在批量上撞见，批量发现的真坑要沉淀回 L0。本报告 B1--B7 属 L1--L3 层读数。布局与职能：

\begin{table}[htbp]
\centering\small
\caption{\texttt{bench/} 现场布局（不熟悉仓库的读者对照此表读结果目录名）}
\label{tab:benchlayout}
\begin{tabular}{@{}lp{10.8cm}@{}}
\toprule
路径 & 内容 \\
\midrule
\texttt{bench/py/} & 37 个评测脚本，分四簇——\textbf{批量驱动}：\texttt{stagerun.py} + \texttt{stage\_\{ingest,parse,xlat,compile,fixloop\}}（五段 DAG，records 逐格 append-only 落盘）；\textbf{七臂 harness}：\texttt{parsebench}、\texttt{compilebench\_v3}、\texttt{xlatbench}、\texttt{qualbench}、\texttt{e2e\_mock\_bench}、\texttt{e2e\_real\_bench}、\texttt{fixloop\_bench}、\texttt{alignbench}、\texttt{validbench}、\texttt{gullet\_bench}、\texttt{fixture\_assert}；\textbf{门禁监控}：\texttt{gate\_scorecard}（M2 出口门）、\texttt{status\_panel}、\texttt{nightwatch}、\texttt{rundiff}、\texttt{triage}、\texttt{wave}、\texttt{stage\_timing}；\textbf{支撑件}：\texttt{benchlib} 公共库、\texttt{translators\_bench}（mock 翻译厂）、\texttt{corpus/} 语料构建管线、\texttt{iclr\_*} 对照语料渠、\texttt{report/} 一次性审计、\texttt{scratch/} 探针 \\
\texttt{bench/ts/} & JS 侧三库横评 harness（latex-utensils/unified-latex/tree-sitter-latex），D0 选型证据产地 \\
\texttt{bench/corpus*} & 语料族：\texttt{corpus} 39 篇手挑陷阱、\texttt{corpus\_v2} 139 篇分层、\texttt{corpus\_v3} 13{,}266 篇八层钉版（\S\ref{sec:corpus}）、\texttt{corpus\_daily} RSS 日更渠（$\sim$1{,}200 篇/日）、\texttt{corpus\_iclr*} ICLR 对照渠 \\
\texttt{bench/fixtures/} & 10 件手构陷阱 \texttt{.tex}（\texttt{@Tnn}/\texttt{@Wnn}/\texttt{@Xn} 断言标记：T 为结构陷阱、W 为野外发现机制、X 为翻译层；逐字节即语义，不做任何格式化） \\
\texttt{bench/results/} & 评测产物区：每次运行一个日期目录（377+ 轮累积），含 records/summary/jsonl/PDF 中间件；大体量不入档，本报告读数全部由其抽取进 \texttt{data/} \\
\texttt{bench/work\_*/} & 编译/fixloop 逐格工作区（解压源树、aux、PDF；重产物 gitignored） \\
\bottomrule
\end{tabular}
\end{table}

\subsection{语料分层构造}
\label{sec:corpus}

corpus\_v3 是钉版批量语料：每个 cell 记录 \texttt{(channel, item, member, blob\_sha256)} 四元组，渠道分 IA \texttt{arxiv-bulk} 月 chunk、e-print 单发、HF \texttt{arxiv-latex-5T}、scholarweave 脱水四种。八层构成见图~\ref{fig:corpus}：核心/补强/扩展/热层为早期层，2026-09-19 评测/开发分轨扩层后收至 13{,}266 篇。

\begin{figure}[htbp]
\centering
\begin{tikzpicture}
\begin{axis}[
  xbar stacked,
  width=13cm, height=4.6cm,
  xmin=0, xmax=14000,
  xlabel={千篇}, ytick=data,
  xtick={0,2000,4000,6000,8000,10000,12000,14000},
  xticklabels={0,2,4,6,8,10,12,14},
  yticklabels={corpus\_v3 八层},
  legend style={at={(0.5,-0.28)},anchor=north,legend columns=4,font=\scriptsize},
  bar width=8mm,
  x label style={font=\small},
]
\addplot+[fill=cA!75,draw=cA] coordinates {(1000,0)};
\addplot+[fill=cB!75,draw=cB] coordinates {(200,0)};
\addplot+[fill=cD!75,draw=cD] coordinates {(3866,0)};
\addplot+[fill=cE!75,draw=cE] coordinates {(166,0)};
\addplot+[fill=cF!75,draw=cF] coordinates {(2000,0)};
\addplot+[fill=cC!75,draw=cC] coordinates {(1514,0)};
\addplot+[fill=cB!40,draw=cB] coordinates {(1500,0)};
\addplot+[fill=black!35,draw=black!60] coordinates {(3020,0)};
\legend{core 1000, booster 200, expand 3866, hot 166, dev\_vol 2000, dev\_recent 1514, dev\_failmine 1500, holdout 3020}
\end{axis}
\end{tikzpicture}
\caption{corpus\_v3 八层构成（13{,}266 篇 / 46GB）。holdout 层为 EVAL\_ONLY 治理——\texttt{dev\_layers()} 实测排除，仅评测可读；dev\_* 三层与 expand 为开发期增量；core 为月度簇分层随机（30 簇，old 200/new 800）。}
\label{fig:corpus}
\end{figure}

治理规则：holdout 层（3{,}020 篇）登记 \texttt{EVAL\_ONLY\_LAYERS}，所有开发期采样经 \texttt{dev\_layers()} 过滤；语料扩展只加层不删层；cell 数据 gitignored，入库物为 manifest + 选择报告 + 机制台账 \texttt{mechanisms.jsonl}（两者快照均随档 \texttt{refs/}；台账是 benchmark $\to$ 归因 $\to$ 规则沉淀的反馈环载体）。

\subsection{门槛体系}

各臂门槛（语料规格 §8，随档 \texttt{refs/}）：parsebench——parse ok 100\%、identity $\geq 99.5\%$、leak $\leq 0.15\%$、dead$\cdot$orphan $=0$、flatten $\geq 99\%$；M2 出口门（stagerun scorecard）——union pdf $\geq 90\%$、clean $\geq 90\%$。统计口径：核心层加权池化率 + 宏平均双报，Wilson + 月簇稳健 bootstrap CI。
        % 评测体系设计（七臂 + 语料 + 门槛）
\section{指标演化史}

本节数据来自三路考古：全部 1{,}648 条历史 commit 逐条扫描（里程碑识别）、评测结果流水 377 个日期目录的读数回填、车队台账与上一期复盘（口径锚点）。本文档自包含：原始数据入库于随档 \texttt{data/} 子目录（scorecard/parse/corpus/rules-tests/qual-xlat 五条时间线 CSV + 377 行评测活动流水含各轮 run 目录名 + 1{,}648 行 commit 全录 + milestones.md），引用源文档时点快照在随档 \texttt{refs/}（台账、复盘、评测器与语料规格、8 库横评报告、机制台账）；口径不一的读数分线绘制并注明样本框。

\subsection{开发节奏：六天 1{,}648 个 commit}

\begin{figure}[htbp]
\centering
\begin{tikzpicture}
\begin{axis}[
  width=12.5cm, height=5cm,
  ybar, bar width=9mm,
  ymin=0,
  symbolic x coords={09-14,09-15,09-16,09-17,09-18,09-19},
  xtick=data, x tick label style={font=\small},
  ylabel={commits/日}, y label style={font=\small},
  nodes near coords, nodes near coords style={font=\scriptsize},
]
\addplot+[fill=cA!70,draw=cA] coordinates {(09-14,9) (09-15,123) (09-16,306) (09-17,624) (09-18,265) (09-19,321)};
\end{axis}
\end{tikzpicture}
\caption{commit 日吞吐（累计 1{,}648）。09-17 峰值 624 对应「波次战」（vendor shim/illegal\_unit/段修多车道并行收割）；全程为多代理车队协作模式（overseer 台账记 roster 峰值 $\sim$23 车道），commit 由 leader 会话统一落盘。}
\label{fig:commits}
\end{figure}

项目节奏呈「日主题」结构（图~\ref{fig:commits}、表~\ref{tab:days}）：D0 立项与八库横评定案、D1 产品骨架（七臂同日落地）、D2 v2 切换 + fixloop 接线（最大单日指标跃升）、D3 波次战收割残面、D4 selfimp 常驻环、D5 语料收官 + 全量综合测试。

\begin{table}[htbp]
\centering\small
\caption{六日大事记（commit 考古，精选）}
\label{tab:days}
\begin{tabular}{@{}llp{9.2cm}@{}}
\toprule
日 & 主题 & 标志事件 \\
\midrule
D0 09-14 & 立项定案 & 仓库初始化；\textbf{8 库横评}（39 篇 + 30 fixtures）——miniscanner 陷阱 32/32、identity 259/259 定案自研路线；技术栈 ADR 冻结 \\
D1 09-15 & 产品骨架 & miniscanner 重写为 \texttt{texlate/latex} 九文件；Mouth+Gullet 展开机；compile/xlat/server/阅读器全套；\textbf{七臂同日落地}；fixloop 25 规则移植；语料 $\to$1{,}339 \\
D2 09-16 & 切换+接线 & \textbf{v2 cutover}（splice 残留 1524$\to$0）；\textbf{fixloop 接线}（pdf $+18.9$pt）；expand+hot 层语料 $\to$5{,}410；loop1 $n$=5124 scorecard 97.89\%；qualbench 诞生 \\
D3 09-17 & 波次战 & vendor 23 件翻 492 格；illegal\_unit 108 闭环；ReDoS 修复 pytest 48min$\to$3:46；xlat P0 封死错配通道；\textbf{realn200 98.5\% 真臂验收} \\
D4 09-18 & selfimp 环 & 常驻改进环 23 车道发车；loop2 97.26\%/86.42\%；\textbf{dollar 族 leak 清零}；qualbase 92.9 建档；corpus\_daily 日更渠启动 \\
D5 09-19 & 收官+全量 & \textbf{corpus\_v3 八层 13{,}266 篇}（dev/eval 分治）；完整性审计全绿；\textbf{12 臂综合测试同日落地}（本报告主体）；loop3 在飞 \\
\bottomrule
\end{tabular}
\end{table}

\subsection{编译健康度}

\begin{figure}[htbp]
\centering
\begin{tikzpicture}
\begin{axis}[
  width=14cm, height=8cm,
  date coordinates in=x,
  xtick={2026-09-15,2026-09-16,2026-09-17,2026-09-18,2026-09-19},
  x tick label style={font=\scriptsize},
  ymin=0, ymax=108,
  ylabel={\%}, y label style={font=\small},
  legend style={at={(0.02,0.97)},anchor=north west,font=\scriptsize,legend columns=2},
  grid=major, grid style={black!15},
]
% scorecard union pdf（stagerun records 重评）
\addplot+[color=cA,mark=*,thick] coordinates {
  (2026-09-16,89.17) (2026-09-17,95.99) (2026-09-17,97.35)
  (2026-09-17,97.89) (2026-09-18,97.26) (2026-09-19,98.75)
};
% scorecard clean
\addplot+[color=cB,mark=*,thick] coordinates {
  (2026-09-16,57.20) (2026-09-17,68.91) (2026-09-17,82.76)
  (2026-09-17,84.99) (2026-09-18,86.42) (2026-09-19,88.75)
};
% real 臂
\addplot+[color=cE,mark=triangle*,mark size=3pt] coordinates {
  (2026-09-17,98.5) (2026-09-19,96.7)
};
\addplot+[color=cE,mark=triangle,mark size=3pt,dashed] coordinates {
  (2026-09-17,86.5) (2026-09-19,86.7)
};
% 裸编译基线
\addplot+[color=black!55,mark=square,only marks] coordinates {
  (2026-09-15,71.7) (2026-09-16,70.6) (2026-09-16,71.3) (2026-09-19,90.2)
};
\addplot+[color=cD,dashed] coordinates {(2026-09-15,90) (2026-09-19,90)};
\legend{union pdf（scorecard）, clean（scorecard）, real 臂 union pdf, real 臂 clean, 裸编译 union pdf, M2 门 90\%}
\end{axis}
\end{tikzpicture}
\caption{编译健康度时间线。scorecard 线为 stagerun records 在代码演进下的重评（loop1 $n$=5124 $\to$ loop2 $n$=5219 $\to$ v3all 样本 $n$=80）；real 臂为真 LLM 翻译全链（realn200 $\to$ e2e-v3b $n$=60 终态）；方块为裸编译基线——注意\textbf{基线自身四天涨 $\sim$19pt}（71.7$\to$90.2）：vendored 宏包资产 + tlmgr usertree 积累使工具链整体成熟，「裸引擎撑死七成」的 09-15 结论已不适用今日地板。}
\label{fig:health}
\end{figure}

三条叙事线（图~\ref{fig:health}）：\textbf{fixloop 是单一最大杠杆}——09-16 接产品链当日 union pdf $+18.9$pt（70.6$\to$89.5，同 n180 复跑）；\textbf{波次战把残面打成个位数}——vendor 波翻 492 格 fail 面、illegal\_unit 108/108 闭环、scorecard 一波 89.17$\to$97.89；\textbf{real 与 mock 打平}——realn200 union pdf 98.5\% 证明评测面非虚高，今日 e2e-v3b 终态 96.7\%（7 个 pipe-fail 全被 fixloop 救回、2 格 oversize 跳过）。未达标字面只剩 M2 clean $\geq90\%$（今 88.75\%，$-1.25$pt）。

\subsection{解析侧：identity 守住，leak 四天内归零再现身}

\begin{figure}[htbp]
\centering
\begin{tikzpicture}
\begin{groupplot}[
  group style={group size=2 by 1,horizontal sep=16mm},
  width=7.4cm, height=5cm,
]
\nextgroupplot[
  ybar, bar width=8mm,
  ymin=0,
  symbolic x coords={09-15 核 1955,09-17 核 1955,09-18 核 1955,09-19 全 28904},
  xtick=data, x tick label style={font=\tiny,rotate=18},
  ylabel={leak 命中数}, y label style={font=\small},
  nodes near coords, nodes near coords style={font=\scriptsize},
]
\addplot+[fill=cB!70,draw=cB] coordinates {(09-15 核 1955,58) (09-17 核 1955,57) (09-18 核 1955,0) (09-19 全 28904,76)};
\nextgroupplot[
  ybar, bar width=8mm,
  ymin=99.9, ymax=100.01,
  symbolic x coords={核心 1955,补强 1388,expand 7229,全量 28904},
  xtick=data, x tick label style={font=\tiny,rotate=18},
  ylabel={identity \%}, y label style={font=\small},
  ytick={99.9,99.95,100},
  nodes near coords, nodes near coords style={font=\scriptsize},
]
\addplot+[fill=cA!70,draw=cA] coordinates {(核心 1955,100) (补强 1388,100) (expand 7229,100) (全量 28904,99.99)};
\end{groupplot}
\end{tikzpicture}
\caption{解析侧演化。左：leak 命中数——dollar 族 58/57 条残留于 09-18 c2-dollar 车道清零（核心层口径），09-19 全量 28{,}904 文件上重现 76 条（1{,}845{,}338 chunk 中 0.004\%，新语料表面新案例，非回归）。右：strict identity 全测量期 $\geq$99.99\%（纵轴截断）。}
\label{fig:parse}
\end{figure}

解析侧四天稳态（图~\ref{fig:parse}）：v2 Gullet+Segmenter cutover 当日 splice 残留占位符 1524$\to$0；leak 残留清一色 dollar 族，09-18 单类收口归零——但\textbf{归零是在 1{,}955 文件核心层上}，今日 8 倍体量全量复测重现 76 条命中，说明该族机制在长尾语料上未穷举，属于「已识别族的覆盖率问题」而非新机制。

\subsection{翻译侧：硬契约与质量评分}

\begin{figure}[htbp]
\centering
\begin{tikzpicture}
\begin{groupplot}[
  group style={group size=2 by 1,horizontal sep=16mm},
  width=7.4cm, height=5cm,
]
\nextgroupplot[
  ybar, bar width=7mm,
  ymin=55, ymax=105,
  symbolic x coords={09-15 smoke 48,09-15 contract 80,09-19 全层 271},
  xtick=data, x tick label style={font=\tiny,rotate=18},
  ylabel={hard\_ok \%}, y label style={font=\small},
  nodes near coords, nodes near coords style={font=\scriptsize},
]
\addplot+[fill=cA!70,draw=cA] coordinates {(09-15 smoke 48,64.6) (09-15 contract 80,98.8) (09-19 全层 271,93.7)};
\nextgroupplot[
  ybar, bar width=7mm,
  ymin=88, ymax=96,
  symbolic x coords={09-18 qualbase 1200,09-18 qualdrift 300,09-19 e2ereal 1937},
  xtick=data, x tick label style={font=\tiny,rotate=18},
  ylabel={judge mean}, y label style={font=\small},
  nodes near coords, nodes near coords style={font=\scriptsize},
]
\addplot+[fill=cD!70,draw=cD] coordinates {(09-18 qualbase 1200,92.9) (09-18 qualdrift 300,92.8) (09-19 e2ereal 1937,94.0)};
\end{groupplot}
\end{tikzpicture}
\caption{翻译侧演化。左：swe-2-medium 硬契约率——09-15 smoke 期网关不稳（http\_err 17）压到 64.6\%，契约批 98.8\%，今日全层抽样 93.7\%（新校验口径 v\_err 27 条计入后更严）；右：ESA judge 均分 92.9$\to$94.0（swe-2-max 评审，contested 8.1\%$\to$5.5\% 下降）。}
\label{fig:xlatqual}
\end{figure}

\subsection{评测资产与代码规模}

\begin{figure}[htbp]
\centering
\begin{tikzpicture}
\begin{groupplot}[
  group style={group size=4 by 1,horizontal sep=8mm},
  width=4.1cm, height=4.8cm,
]
\nextgroupplot[
  ybar, bar width=3.2mm, ymin=0,
  symbolic x coords={15,16,17,18,19},
  xtick=data, x tick label style={font=\tiny},
  title={pytest 用例}, title style={font=\tiny},
  nodes near coords, nodes near coords style={font=\tiny,rotate=90},
]
\addplot+[fill=cA!70,draw=cA] coordinates {(15,798) (16,1920) (17,4333) (18,5261) (19,6450)};
\nextgroupplot[
  ybar, bar width=3.2mm, ymin=0,
  symbolic x coords={15,16,17,18,19},
  xtick=data, x tick label style={font=\tiny},
  title={fixloop 规则}, title style={font=\tiny},
  nodes near coords, nodes near coords style={font=\tiny,rotate=90},
]
\addplot+[fill=cB!70,draw=cB] coordinates {(15,31) (16,50) (17,71) (18,106) (19,143)};
\nextgroupplot[
  ybar, bar width=3.2mm, ymin=0,
  symbolic x coords={15,16,17,18,19},
  xtick=data, x tick label style={font=\tiny},
  title={src 文件数}, title style={font=\tiny},
  nodes near coords, nodes near coords style={font=\tiny,rotate=90},
]
\addplot+[fill=cE!70,draw=cE] coordinates {(15,72) (16,89) (17,174) (18,355) (19,425)};
\nextgroupplot[
  ybar, bar width=3.2mm, ymin=0,
  symbolic x coords={c39,v2,+core,+boo,+hot,+exp,+dev4},
  xtick=data, x tick label style={font=\tiny,rotate=28},
  title={语料篇数 /千篇（v3 分层累加）}, title style={font=\tiny},
  nodes near coords, nodes near coords style={font=\tiny,rotate=90},
]
\addplot+[fill=cD!70,draw=cD] coordinates {(c39,0.039) (v2,0.139) (+core,1.0) (+boo,1.2) (+hot,1.366) (+exp,5.232) (+dev4,13.266)};
\end{groupplot}
\end{tikzpicture}
\caption{资产增长（git 历史抽样，同日口径一致）：pytest 798$\to$6{,}450、fixloop 规则 31$\to$143（16 分片）、src 72$\to$425 文件、语料 corpus39/v2 $\to$ v3 八层 13{,}266（另有 corpus\_daily 日更渠 $\sim$1{,}200 篇/日在库外增长）。}
\label{fig:assets}
\end{figure}

\subsection{归因表：哪些改动搬动了哪些指标}

\begin{table}[htbp]
\centering\small
\caption{改动 $\to$ 指标位移归因（commit 级，摘自台账与各波报告）}
\label{tab:attrib}
\begin{tabular}{@{}llp{8.4cm}@{}}
\toprule
日期 & 改动（commit） & 指标位移 \\
\midrule
09-14 & 8 库横评 + miniscanner spike & 主解析器路线定案（陷阱 32/32，identity 259/259） \\
09-15 & corpus\_v3 core+booster + 双引擎基线 & 基线口径建立（union pdf 71.7\%/clean 39.4\%，n180） \\
09-15 & fixloop yaml 引擎落地（\texttt{179f7e5b}） & 25 条规则移植，修复臂诞生 \\
09-16 & fixloop 接产品链（\texttt{3d5de8f7}） & union pdf 70.6$\to$89.5\%（$+18.9$pt，同口径复跑） \\
09-16 & v2 Gullet+Segmenter 默认切换（\texttt{f4616838}） & splice 残留占位符 1524$\to$0（cutover 硬门） \\
09-16 & argspec.json 入包（\texttt{401e9bbf}） & 1820 条宏参规格零装载 GAP 闭合 \\
09-16 & perf-tail 五修 & latex parse 长尾 $-48\%$（51476$\to$26723ms），identity 零回退 \\
09-16 & stagerun loop1 全 DAG（$n$=5124） & scorecard 97.89\%/84.99\%；fixloop rescue 84.0\% \\
09-17 & vendored\_fetch + vendor 23 件（\texttt{a90978ab}） & 绝版宏包缺件面 492 格 fail$\to$73.2\% clean，clean +101 格 \\
09-17 & illegal\_unit 段修波（\texttt{d2c377bd}） & 108/108 闭环零回归，clean 84.76$\to$85.03\% \\
09-17 & ReDoS 原子组修复（\texttt{746237fd}） & held6 卡死 30min+$\to$0.1s；pytest 全量 48min$\to$3:46 \\
09-17 & xlat 非锚定 [n] parse 撤除（\texttt{164a9e0f}，P0） & 译文静默错配通道封死，同批核销 4 P1 \\
09-17 & realn200 真臂验收（\texttt{f2847c33}） & union pdf 197/200=98.5\%，与 mock 打平 \\
09-18 & c2-dollar 车道收口 & leak 57$\to$0/136{,}049 chunks（核心层口径单类清零） \\
09-18 & gate\_scorecard schema v3 & loop1 复跑 97.89/84.99 分毫不差——评测器自证稳定 \\
09-18 & qualbase esa2（$n$=1200） & 质量基线建档 mean 92.9 [92.4,93.3]，contested 8.1\% \\
09-18 & loop2 冻结快照（$n$=5219） & scorecard 97.26\%/86.42\%（clean +1.43pt） \\
09-19 & corpus 8 层扩层（\texttt{13c603cf}） & 语料 5232$\to$13{,}266，EVAL\_ONLY 治理落地 \\
09-19 & 全量综合测试（本报告） & 12 臂同日落地；基线地板升至 90.2\%；stagerun 98.75\%/88.75\% \\
\bottomrule
\end{tabular}
\end{table}
      % 指标演化史（时间线图 + 归因表 + 大事记）
\section{2026-09-19 全量画像}

本节为 corpus\_v3 扩充收官（13{,}266 篇）后的首轮全量综合测试，12 个评测臂同日落地。总表见表~\ref{tab:snapshot}，逐节展开关键分布。

\begin{table}[htbp]
\centering\scriptsize
\caption{2026-09-19 综合测试逐臂指标总表}
\label{tab:snapshot}
\begin{tabular}{@{}llp{7.6cm}p{3.6cm}@{}}
\toprule
臂 & 规模 & 核心指标 & 门槛对比 \\
\midrule
完整性审计 & 13{,}266 cells & sha256 13{,}266/13{,}266，0 dup，0 missing；13 stub 为 B07 有意构造 & EVAL\_ONLY 生效 \\
parsebench & 13{,}253 篇 / 28{,}904 .tex & parse ok \textbf{100.0\%}；identity \textbf{99.99\%}；leak \textbf{0.004\%}（76/1{,}845{,}338）；expand footprint strict 28105/norm 259/div 539；flatten 93.7\% & identity/leak PASS；\textbf{dead ph=3 FAIL}；\textbf{flatten FAIL} \\
compilebench base & 500 格 $\times$2 & xelatex 56.3\% clean/80.3\% pdf；tectonic 38.6\%/62.4\%；联合 pdf \textbf{90.4\%} & --- \\
compilebench zh & 500 格 & xelatex \textbf{72.4\%/85.4\%}（+16pt）；tectonic 45.7\%/59.8\% & --- \\
fixloop（v2 40） & 80 格 & xelatex 34 FAIL$\to$\textbf{25 救回 74\%}；tectonic 1/16（eps\_route 主导） & --- \\
alignbench & 812 对 & 门全过；765/812 对 named-dests$=0$，真实覆盖 42 对（retention p50 1.0） & 覆盖不足 \\
e2e\_mock & corpus39 & pipe-xel \textbf{33/39 clean} vs base-xel 19/39；引入回归 3 & --- \\
e2e\_real & n=24+60 & 58/58 翻译；chunk ok 6053/7580；splice 残留 \textbf{0}；pipe-xel clean 43/fail 7；fixloop 再救 13；\textbf{0 引入回归} & --- \\
xlatbench & 271 样本 & hard\_ok \textbf{93.7\%}；http\_err 0；延迟 p50 4.2s/p95 29.7s & --- \\
qualbench & 324 篇/1937 chunk & judge 均分 \textbf{94.0}，median 95；$\geq$90 占 90.6\%；contested 5.5\%；critical 1 & --- \\
gullet & 200 文档 & 199/200 展开成功；median 39.6ms/19 steps & --- \\
stagerun DAG & 80 篇跨层 & \textbf{union pdf 79/80=98.75\%}；clean 71/80=88.75\% & M2 pdf 门 \textbf{PASS} \\
\bottomrule
\end{tabular}
\end{table}

\subsection{全链漏斗：stagerun DAG}

五阶段批量驱动对 80 篇跨层样本的实测漏斗（图~\ref{fig:funnel}）：ingest 80 $\to$ parse 79 ok（1 reject）$\to$ xlat mock 臂 76 ok + 3 partial $\to$ compile zh 臂 61 clean + 13 partial + 5 fail + 1 skip $\to$ fixloop 在 21 格上动作（fail$\to$clean 5、partial$\to$clean 5、partial 保持 8、clean 复核 3）。union 终态 \textbf{clean 71 + partial 8 = pdf 79/80（98.75\%）}，过 M2 union-pdf $\geq90\%$ 门；clean 88.75\% 距 90\% 门差 1.25pt。同格对照臂 base（裸编译 50 clean + 19 partial + 10 fail）给出归因锚点：管线+fixloop 相对裸编译 \textbf{pdf +10 格（69$\to$79）、clean +21 格（50$\to$71）}。

\begin{figure}[htbp]
\centering
\begin{tikzpicture}
\begin{groupplot}[
  group style={group size=2 by 1,horizontal sep=14mm},
]
\nextgroupplot[
  xbar,
  width=8.6cm, height=4.6cm,
  xmin=0, xmax=90,
  xlabel={格数}, ytick=data,
  yticklabels={union pdf（71 clean+8 partial）,compile zh pdf（61+13）,xlat ok,parse ok,ingest},
  y tick label style={font=\scriptsize},
  nodes near coords, nodes near coords style={font=\scriptsize},
  bar width=4.2mm,
  x label style={font=\small},
]
\addplot+[fill=cA!70,draw=cA] coordinates {(79,0) (74,1) (76,2) (79,3) (80,4)};
\nextgroupplot[
  ybar, bar width=5.5mm,
  width=6.4cm, height=4.6cm,
  ymin=0, ymax=90,
  symbolic x coords={base 裸编译,zh+fixloop union},
  xtick=data, x tick label style={font=\scriptsize},
  ylabel={格数}, y label style={font=\small},
  legend style={at={(0.5,-0.24)},anchor=north,legend columns=2,font=\scriptsize},
  nodes near coords, nodes near coords style={font=\scriptsize},
]
\addplot+[fill=cD!70,draw=cD] coordinates {(base 裸编译,50) (zh+fixloop union,71)};
\addplot+[fill=cC!75,draw=cC] coordinates {(base 裸编译,19) (zh+fixloop union,8)};
\legend{clean,partial（pdf\~{}）}
\end{groupplot}
\end{tikzpicture}
\caption{stagerun 五阶段漏斗与同格对照（80 篇跨层样本，xlat 为 mock 上游）。左：主流水线 zh 臂逐段存活数；右：同一 80 格上 base 裸编译 vs zh+fixloop union 终态——管线相对裸编译 pdf +10 格、clean +21 格（fixloop 救回 fail 5 格全升 clean、partial 5 格升 clean）。}
\label{fig:funnel}
\end{figure}

\subsection{翻译质量：qualbench ESA 评分}

swe-2-medium 译文对 324 篇论文 1{,}937 个 chunk 的 LLM-judge 评分（judge=swe-2-max，二裁 swe-2-high，contested 5.5\% 触发二裁）：mean 94.0 / median 95，$\geq$90 分占 90.6\%，critical 级缺陷仅 1 例（non-translation）。分数段与文类分布见图~\ref{fig:qual}：section\_title 最高（96.9），caption 93.2，para 92.2——长散文段仍是质量洼地。

\begin{figure}[htbp]
\centering
\begin{tikzpicture}
\begin{groupplot}[
  group style={group size=2 by 1,horizontal sep=14mm},
  width=7.2cm, height=5cm,
]
\nextgroupplot[
  ybar, bar width=7mm,
  ymin=0,
  symbolic x coords={$\geq$90,75--89,55--74,$<$55},
  xtick=data, x tick label style={font=\scriptsize},
  ylabel={chunk 数}, y label style={font=\small},
  nodes near coords, nodes near coords style={font=\scriptsize},
]
\addplot+[fill=cA!70,draw=cA] coordinates {($\geq$90,1756) (75--89,133) (55--74,41) ($<$55,7)};
\nextgroupplot[
  ybar, bar width=5mm,
  ymin=88, ymax=98,
  symbolic x coords={para,caption,section\_title},
  xtick=data, x tick label style={font=\scriptsize,rotate=15},
  ylabel={mean 分}, y label style={font=\small},
  nodes near coords, nodes near coords style={font=\scriptsize},
]
\addplot+[fill=cD!70,draw=cD] coordinates {(para,92.2) (caption,93.2) (section\_title,96.9)};
\end{groupplot}
\end{tikzpicture}
\caption{qualbench 评分分布（左：分数段；右：per-kind 均分）。1{,}756/1{,}937 chunk $\geq$90；para 类是主要拉低项（flag 榜：term\_inconsistency 229、grammar 197、mistranslation 132、hallucinated\_content 66）。}
\label{fig:qual}
\end{figure}

\subsection{编译臂：baseline vs zh 条件}

500 格抽样双引擎对照（图~\ref{fig:compile}）：zh 条件（产品链 normalize + ctex 注入后编译）反而比裸编译高 16pt clean——normalize 顺带修复了源文件的编码/格式缺陷。tectonic 在 zh 臂 clean 升而 pdf 略降，EPS 图源是其死穴（见 \S\ref{sec:failures}）。

\begin{figure}[htbp]
\centering
\begin{tikzpicture}
\begin{axis}[
  ybar, bar width=6mm,
  width=11.5cm, height=5.4cm,
  ymin=0, ymax=100,
  ylabel={\%}, y label style={font=\small},
  symbolic x coords={xelatex clean,xelatex pdf,tectonic clean,tectonic pdf,union pdf},
  xtick=data, x tick label style={font=\scriptsize},
  legend style={at={(0.5,-0.22)},anchor=north,legend columns=2,font=\scriptsize},
  nodes near coords, nodes near coords style={font=\scriptsize},
]
\addplot+[fill=black!45,draw=black!60] coordinates {(xelatex clean,56.3) (xelatex pdf,80.3) (tectonic clean,38.6) (tectonic pdf,62.4) (union pdf,90.4)};
\addplot+[fill=cA!70,draw=cA] coordinates {(xelatex clean,72.4) (xelatex pdf,85.4) (tectonic clean,45.7) (tectonic pdf,59.8) (union pdf,0)};
\legend{baseline（裸编译）, zh 条件（normalize+ctex 后）}
\end{axis}
\end{tikzpicture}
\caption{compilebench 500 格：baseline vs zh 条件。zh 臂 xelatex clean +16.1pt——normalize 修复了源级缺陷；union pdf 90.4\%（双引擎取优）。}
\label{fig:compile}
\end{figure}

\subsection{翻译硬契约：xlatbench}

271 个抽样上下文对 swe-2-medium 的硬契约测试：hard\_ok 254/271=93.7\%，http\_err 0。延迟 p50 4.2s / p90 16.2s / p95 29.7s / max 104.4s——bg 批量档 in-gate 排队在长尾可见。token 经济：tok\_in 77.4k / tok\_out 40.8k（均 285/150 per chunk），reasoning 均 93 tok，prompt 缓存命中 0\%（抽样源分散）。per-kind 契约率：stress 4/4、para 64/67、caption 62/67、item 61/66、footnote 58/60——footnote/item 类是短板。失败类：cs\_dropped$\times$8、ph\_miss$\times$5、validator$\times$3、ord\_hard$\times$1；ord\_soft 78 条软序告警（不改契约，但提示换序倾向）。

\subsection{e2e 全链：mock 与 real 双臂}

e2e\_mock（corpus39）：pipe-xel 33/39 clean vs base-xel 19/39——管线净正，翻译+splice 步骤不损反而增益（normalize 修复）；3 格引入回归。e2e\_real n=60：58/58 翻译完成，chunk ok 6053/7580（partial 43/fault 8），splice 残留占位符 0（PASS），pipe-xel clean 43/fail 7/partial 6，fixloop 再救 13 格，\textbf{0 管线引入回归}。编译失败类：missing\_file$\times$7、missing\_graphic$\times$3、killed\_by\_signal$\times$3、driver\_fatal$\times$3——以语料源缺资产为主。
     % 2026-09-19 全量画像（逐臂 + 漏斗 + 分布）
\section{失败归因图谱}
\label{sec:failures}

\subsection{编译侧失败类}

e2e\_real n=60 的编译失败归因（图~\ref{fig:failclass} 左）：missing\_file$\times$7 与 missing\_graphic$\times$3 合计过半——语料源自身缺资产（arXiv e-print 打包不全），非管线引入；killed\_by\_signal$\times$3 为超时/资源类；driver\_fatal$\times$3 为引擎级崩溃。

fixloop 救回矩阵（corpus\_v2 40 篇对照）：xelatex 34 个 FAIL 格救回 25（74\%），主力规则为 \texttt{install\_file}（14 触发，8 格出 pdf）、\texttt{install\_tfm}（12 触发，7 格）、\texttt{static\_precheck}（40 格全触发）；装包榜首位 subfigure.sty$\times$10、revtex4.cls$\times$9。tectonic 臂仅救 1/16——\textbf{eps\_route reject 是结构性死穴}（11 格全拒）：EPS 图源在 tectonic/xdvipdfmx 通路无解，唯一出路是路由 xelatex。另有 latex209\_reject 闸拦下 LaTeX2.09 古文 8 格（该类论文需 209$\to$2e 重写通道，非装包可救）。

\subsection{翻译侧失败类}

xlatbench 失败类分布（图~\ref{fig:failclass} 右，17 个 hard\_fail 全归因）：\textbf{cs\_dropped$\times$8 居首}——控制序列在译文中丢失（占位符协议被破坏的最常见形态）；ph\_miss$\times$5（占位符未还原）、validator$\times$3、ord\_hard$\times$1；另 ord\_soft 78 条为不破契约的 chunk 顺序软告警。qualbench flag 榜（1{,}937 judged chunk）：term\_inconsistency 229、grammar 197、mistranslation 132、fluency\_register 99、hallucinated\_content 66、untranslated\_spans 37、over\_translation 23、placeholder\_broken 8。

\begin{figure}[htbp]
\centering
\begin{tikzpicture}
\begin{groupplot}[
  group style={group size=2 by 1,horizontal sep=16mm},
  width=7.4cm, height=5.2cm,
]
\nextgroupplot[
  xbar, bar width=3.6mm,
  xmin=0,
  symbolic y coords={driver\_fatal,killed\_by\_signal,missing\_graphic,missing\_file},
  ytick=data, y tick label style={font=\scriptsize},
  xlabel={格数}, x label style={font=\small},
  nodes near coords, nodes near coords style={font=\scriptsize},
]
\addplot+[fill=cB!70,draw=cB] coordinates {(3,driver\_fatal) (3,killed\_by\_signal) (3,missing\_graphic) (7,missing\_file)};
\nextgroupplot[
  xbar, bar width=3.6mm,
  xmin=0,
  symbolic y coords={ord\_soft,ord\_hard,validator,ph\_miss,cs\_dropped},
  ytick=data, y tick label style={font=\scriptsize},
  xlabel={样本数}, x label style={font=\small},
  nodes near coords, nodes near coords style={font=\scriptsize},
]
\addplot+[fill=cC!75,draw=cC] coordinates {(78,ord\_soft) (1,ord\_hard) (3,validator) (5,ph\_miss) (8,cs\_dropped)};
\end{groupplot}
\end{tikzpicture}
\caption{失败类分布。左：e2e\_real n=60 编译失败归因——语料源缺资产主导；右：xlatbench 271 样本翻译失败/告警类——cs\_dropped 为头号硬失败。}
\label{fig:failclass}
\end{figure}

\subsection{门槛 FAIL 与已知缺口}

\begin{enumerate}
  \item \textbf{dead protect ph=3}（parsebench）：死占位符残留 3 例，门为 0——唯一 gate FAIL 项，值小但需逐例归因。
  \item \textbf{flatten 93.7\% $<99\%$}：已知 errata，orphan tex 清单在 papers.json（多层语料中 \texttt{.tex} 不进主文档 flatten 面的占比）。
  \item \textbf{alignbench 覆盖薄}：812 对中 765 对 named-dests 为零——保留率指标实际只在 42 个 hyperref 对上有意义，指标口径需复核（分母该是所有对还是仅 hyperref 对）。
  \item \textbf{tectonic 短板}：eps\_route 无解 + biblatex$\leftrightarrow$biber 版本错配——EPS/bib 论文在 tectonic 下是先天缺省，路由层应前置判走 xelatex。
  \item \textbf{cs\_dropped 归因待定}：prompt 侧约束不足 vs 网关侧截断，需在 xlatbench 加对照实验分辨。
  \item \textbf{stagerun scorecard 口径}：records 混存首轮欠采的 stale skip 行会污染读数（本批 989 条把 union pdf 冲成 7.4\% 假象；按 ingest 样本口径重算 98.75\%）——工具层待加 \texttt{--since}/样本过滤参数。
  \item \textbf{未测臂}：translators\_bench、validbench（B6）、nightwatch、rundiff、triage、stage\_timing、wave、status\_panel——监控/辅助工具为主，本轮未覆盖。
\end{enumerate}
     % 失败归因图谱 + 已知缺口
\appendix
\section{口径与数据说明}

\subsection{统计口径纪律}

\begin{itemize}
  \item 每个数字标注规模 $n$ 与样本框；不同 $n$/口径的读数不连成同一演化线（图注注明口径切换点）。
  \item scorecard 读数一律按 ingest 样本集过滤后计算——stagerun records 为 append-only，跨轮重跑会残留 stale 行，原始行计数会把未参与样本的 skip 计入分母。
  \item clean/pdf 为编译终态分级：clean = 无 warning 出 PDF；pdf\~{} = 带可接受瑕疵出 PDF（acceptable/best\_effort/dirty 三档）；union pdf = clean + pdf\~{}。
  \item 延迟为网关端到端秒；bg 批量档 in-gate 排队 $\sim$120s，拥塞期 p95 会被队列时间放大，读数与空闲档不可比。
  \item qualbench 评分由 LLM judge（swe-2-max）产出，contested 判例经第二 judge（swe-2-high）仲裁；judge 自身偏差未做人工标定，分数应读作相对分布而非绝对正确率。
\end{itemize}

\subsection{结果目录索引（2026-09-19 批）}

\begin{center}\small
\begin{tabular}{@{}ll@{}}
\toprule
臂 & \texttt{bench/results/} 目录 \\
\midrule
parsebench 全量 & \texttt{parsebench-v3-all-2026-09-19/} \\
compilebench baseline & \texttt{compilebench-v3-2026-09-19/} \\
compilebench zh & \texttt{compilebench-v3-zh-2026-09-19/} \\
xlatbench 全层 & \texttt{xlatbench-v3all-2026-09-19/} \\
qualbench & \texttt{qualbench-e2ereal-2026-09-19/report.md} \\
e2e\_real & \texttt{e2e-real-v3-2026-09-19-*/}（n=24）、\texttt{e2e-real-v3b-*（n=60）} \\
e2e\_mock & \texttt{e2e-mock-2026-09-19-2026-09-19/} \\
alignbench & \texttt{alignbench-e2ereal-2026-09-19/} \\
gullet & \texttt{gullet-v3-2026-09-19/} \\
stagerun DAG & \texttt{stagerun-v3all-2026-09-19/records/} \\
fixloop & \texttt{fixloop-corpusv2-2026-09-15/} \\
\bottomrule
\end{tabular}
\end{center}

配套文档：本批 Markdown 版总表 \texttt{docs/research/2026-09-19-bench-metrics-corpusv3.md}；上一期指标复盘 \texttt{docs/research/review-2026-09-18/README.md}（含 metrics-timeline.svg）；评测器规格 \texttt{docs/10-benchmark-suite.md}；语料规格 \texttt{docs/09-benchmark-corpus.md}；机制台账 \texttt{bench/corpus\_v3/mechanisms.jsonl}。

\subsection{过程事故记录（评测工具链教训）}

本批测试记录的 harness 层教训，供后续批次回避：xlatbench 的 \texttt{manifest.parent} 必须等于语料根（跨目录 manifest 会 cell 解析失败）；\texttt{--per-kind} 是全局限额非分层配额；pgrep 监控须用 \texttt{[x]pattern} 括号形防自匹配；compilebench \texttt{--gen-sample} 为分段设计（采样后即返回，编译需无参二次调用）；e2e 同 \texttt{--tag} 冻结样本，重抽样须换 tag+seed；gullet \texttt{--out} 续跑保留 stale err 行须换新目录；stagerun 下游阶段须显式传 \texttt{--ids} 防跨层欠采；多会话共仓 commit 须走私有 \texttt{GIT\_INDEX\_FILE}（共享 index 有 sweep/reset 竞争窗）。
     % 口径说明 + 结果目录索引

\end{document}
